% The showcase: the family as a reader or a designer meets it first. A cover,
% the alphabets, sizes, the OpenType features each shown off and on, the
% scripts, and two book pages. samples.tex is the diagnostic specimen
% (coverage, the four styles line by line), textpages.tex the long reading
% one; this one is for looking at.
%   lualatex --output-directory=build/proofs proofs/specimen.tex   (twice; or make proof)
\documentclass[11pt]{article}
\usepackage[a4paper,top=22mm,bottom=26mm,inner=24mm,outer=24mm]{geometry}
\usepackage{fontspec}
\usepackage{needspace}
\usepackage{fancyhdr}
\usepackage{multicol}
\usepackage{array}
\usepackage[hidelinks]{hyperref}
\defaultfontfeatures{Path=../fonts/otf/, Renderer=HarfBuzz}
\setmainfont{MyiaDidot-Regular.otf}[BoldFont=MyiaDidot-Bold.otf, ItalicFont=MyiaDidot-Italic.otf,
  BoldItalicFont=MyiaDidot-BoldItalic.otf]
\newfontfamily\R{MyiaDidot-Regular.otf}
\newfontfamily\I{MyiaDidot-Italic.otf}
\newfontfamily\B{MyiaDidot-Bold.otf}
\newfontfamily\BI{MyiaDidot-BoldItalic.otf}
\newcommand\smc[1]{{\addfontfeature{Letters=SmallCaps}#1}}
\newcommand\allsc[1]{{\addfontfeature{Letters={SmallCaps,UppercaseSmallCaps}}#1}}
\newcommand\feat[2]{{\addfontfeature{RawFeature=#1}#2}}
\newcommand\lang[2]{{\addfontfeature{Language=#1}#2}}
\newcommand\cyr[2]{{\addfontfeature{Script=Cyrillic,Language=#1}#2}}
\newfontlanguage{IPA}{IPPH}
\newcommand\sz[2]{\fontsize{#1}{#2}\selectfont}
\InputIfFileExists{../build/proofs/buildinfo.tex}{}{\def\buildinfo{unversioned build}}
\pagestyle{fancy}\fancyhf{}\renewcommand\headrulewidth{0pt}
\fancyhead[R]{\footnotesize\allsc{\leftmark}}
\fancyfoot[L]{\footnotesize Myia Didot · specimen · \buildinfo}
\fancyfoot[R]{\footnotesize\thepage}
% a part of the specimen: a new page, its name in the running head and the contents
\newcommand\spart[1]{\clearpage\markboth{#1}{#1}\phantomsection\addcontentsline{toc}{section}{#1}%
  {\sz{9}{11}\allsc{Myia Didot}}\par\vspace{2pt}{\sz{26}{30}#1}\par\vspace{18pt}}
\newcommand\sub[1]{\par\needspace{5\baselineskip}\vspace{12pt}{\sz{9}{11}\allsc{#1}}\par\vspace{4pt}\nopagebreak}
\setlength{\parindent}{0pt}
\setlength{\parskip}{0pt}
\linespread{1.1}
\begin{document}

% -- cover --------------------------------------------------------------------
\thispagestyle{empty}
\vspace*{30mm}
{\sz{96}{96}Myia Didot}\par\vspace{10mm}
{\sz{20}{26}A Didot for text, in four styles, with real small capitals,\\
for Latin, Greek and Cyrillic}\par\vspace{18mm}
{\sz{30}{40}\R Regular\par\I Italic\par\B Bold\par\BI Bold Italic\par}
\vfill
{\sz{12}{16}\allsc{Aa Bb Cc Dd Ee Ff Gg Hh Ii Jj}\quad Ἀα Ββ Γγ Δδ\quad Жж Щщ Яя\par}
\vspace{6mm}
{\footnotesize A fork of Theano Didot (Alexey Kryukov), its italics after Old Standard;
SIL Open Font License 1.1.\par \buildinfo\par}
\clearpage

\thispagestyle{fancy}\markboth{Contents}{Contents}
{\sz{26}{30}Contents}\par\vspace{12pt}
\begingroup\renewcommand\contentsname{}\vspace*{-36pt}\tableofcontents\endgroup

% -- the four styles at display size --------------------------------------------
\spart{The four styles}
\newcommand\four[1]{{\R #1}\par{\I #1}\par{\B #1}\par{\BI #1}\par}
{\sz{40}{46}\four{Hamburgefonstiv}}
\vspace{10pt}
{\sz{17}{21}\four{ABCDEFGHIJKLMNOPQRSTUVWXYZ}}
\vspace{6pt}
{\sz{17}{21}\four{abcdefghijklmnopqrstuvwxyz ß fi fl}}
\vspace{6pt}
{\sz{17}{21}\four{\allsc{Small capitals ÆŒÐÞ} 0123456789 {\addfontfeature{Numbers=OldStyle}0123456789}}}
\vspace{6pt}
{\sz{17}{21}\four{\&?!¿¡«»“”‘’()[]\{\}§¶†‡€\$£¥\%‰}}

% -- sizes ----------------------------------------------------------------------
\spart{From display to footnote}
\newcommand\wf[3]{{\sz{#1}{#2}#3}\par\vspace{2pt}}
\wf{72}{78}{Didot}
\wf{48}{54}{Imprimerie royale}
\wf{36}{42}{\I Le point typographique}
\wf{28}{33}{Firmin Didot, graveur, 1784}
\wf{20}{24}{\B The contrast between thick and thin strokes}
\wf{16}{20}{Ὅμηρος, Ἰλιάς · Война и мир · Łódź · Ærøskøbing · Hồ Chí Minh}
\wf{14}{17}{It is a truth universally acknowledged, that a single man in possession of a good fortune, must be in want of a wife.}
\wf{12}{15}{It is a truth universally acknowledged, that a single man in possession of a good fortune, must be in want of a wife.}
\wf{10.5}{13.5}{It is a truth universally acknowledged, that a single man in possession of a good fortune, must be in want of a wife. However little known the feelings or views of such a man may be on his first entering a neighbourhood, this truth is so well fixed in the minds of the surrounding families.}
\wf{9}{11.5}{It is a truth universally acknowledged, that a single man in possession of a good fortune, must be in want of a wife. However little known the feelings or views of such a man may be on his first entering a neighbourhood, this truth is so well fixed in the minds of the surrounding families, that he is considered the rightful property of some one or other of their daughters.}
\wf{7}{9}{It is a truth universally acknowledged, that a single man in possession of a good fortune, must be in want of a wife. However little known the feelings or views of such a man may be on his first entering a neighbourhood, this truth is so well fixed in the minds of the surrounding families, that he is considered the rightful property of some one or other of their daughters.}

% -- features -------------------------------------------------------------------
\spart{OpenType features}
{\sz{10}{13}Each line: the text as typed, then with the feature on.}\par\vspace{10pt}
\newcommand\frow[3]{\texttt{\footnotesize #1} & #2 & #3\\[5pt]}
\begingroup\sz{14}{18}
\begin{tabular}{@{}>{\raggedright}p{14mm}p{68mm}p{70mm}@{}}
\frow{smcp}{The Didot family in Paris}{\smc{The Didot family in Paris}}
\frow{c2sc}{UNESCO and NATO}{\allsc{UNESCO and NATO}}
\frow{onum}{in 1784, 1798 and 1801}{\feat{+onum}{in 1784, 1798 and 1801}}
\frow{pnum}{1111 · 1789 · 2026}{\feat{+pnum}{1111 · 1789 · 2026}}
\frow{frac}{1/2 3/4 22/7}{\feat{+frac}{1/2 3/4 22/7}}
\frow{sups}{E = mc2, H2O, x3}{E = mc\feat{+sups}{2}, H\feat{+subs}{2}O, x\feat{+sups}{3}}
\frow{zero}{0 1 0 1}{\feat{+zero}{0 1 0 1}}
\frow{case}{¿(¡HOLA!)? «ABC» – [X]}{\feat{+case}{¿(¡HOLA!)? «ABC» – [X]}}
\frow{liga}{\feat{-liga}{fi fl ff ffi ffl fb fk}}{fi fl ff ffi ffl fb fk}
\frow{dlig}{act aspect Quand The}{\feat{+dlig}{act aspect Quand The}}
\frow{hist}{Kunst ist lang}{\feat{+hist}{Kunst ist lang}}
\frow{calt}{\feat{-calt}{fît naïf fľaša}}{fît naïf fľaša}
\frow{ss01}{Ş ş Ţ ţ}{\feat{+ss01}{Ş ş Ţ ţ}}
\frow{ss02}{\allsc{STRAẞE, Maß}}{\feat{+ss02}{\allsc{STRAẞE, Maß}}}
\frow{locl}{{\I\cyr{Russian}{бгдпт}} (Russian)}{{\I\cyr{Serbian}{бгдпт}} (Serbian)}
\frow{locl}{fi Íjssel (English)}{\lang{Turkish}{fi} (Turkish) \lang{Dutch}{ÍJssel} (Dutch)}
\frow{cpsp}{\feat{-cpsp}{TYPOGRAPHIE}}{TYPOGRAPHIE}
\end{tabular}
\endgroup

% -- scripts --------------------------------------------------------------------
\spart{Latin, Greek, Cyrillic}
\sub{Latin, with Vietnamese and African orthographies}
{\sz{13}{17}Tous les êtres humains naissent libres et égaux en dignité et en droits.
Wszyscy ludzie rodzą się wolni i równi. Árvíztűrő tükörfúrógép. Pijamalı hasta
yağız şoföre çabucak güvendi. Tất cả mọi người sinh ra đều được tự do và bình đẳng
về nhân phẩm và quyền. Hausa: ƙasa, ɗaki, ɓera; Ewe: Eʋegbe, ɖevi, ƒe, ɛ, ɔ.\par}
\sub{Greek, monotonic and polytonic}
{\sz{13}{17}Ἐν ἀρχῇ ἦν ὁ λόγος, καὶ ὁ λόγος ἦν πρὸς τὸν θεόν. Ἄνδρα μοι ἔννεπε,
Μοῦσα, πολύτροπον, ὃς μάλα πολλὰ πλάγχθη. Όλοι οι άνθρωποι γεννιούνται ελεύθεροι
και ίσοι στην αξιοπρέπεια και τα δικαιώματα. {\I Ἡ Ὀδύσσεια τοῦ Ὁμήρου.}\par}
\sub{Cyrillic, with Bulgarian and Serbian forms}
{\sz{13}{17}Все люди рождаются свободными и равными в своём достоинстве и правах.
Всі люди народжуються вільними. {\I Съешь же ещё этих мягких французских булок.}
\cyr{Bulgarian}{Всички хора се раждат свободни} · {\I\cyr{Serbian}{Сва људска бића рађају се слободна.}}\par}
\sub{Phonetics}
{\sz{13}{17}\lang{IPA}{ðə ˈnɔɹθ ˌwɪnd ən ðə ˈsʌn wɚ dɪˈspjutɪŋ ˈwɪtʃ wəz ðə ˈstɹɔŋɡɚ ·
ɓ ɗ ʄ ɠ ʛ ɖ ɭ ɳ ɽ ʈ ʂ ʐ ɻ ʃ ʒ ɕ ʑ ɸ β θ ð χ ʁ ħ ʕ ʔ ǀ ǁ ǂ ǃ ʘ}\par}

% -- book pages -----------------------------------------------------------------
\spart{A book page}
\begin{multicols}{2}
\sz{10.5}{14}\setlength{\parindent}{1.2em}\addfontfeature{Numbers=OldStyle}
\smc{When the Didot family} set up their presses in Paris in the last decades
of the eighteenth century, the printed page in France had already been through
two hundred years of quiet change, and almost nobody outside the trade noticed
what happened next. François-Ambroise Didot, the father, was a printer before he
was a designer of type, and his first concern was not the shape of the letters
but the way they were measured: the \emph{point Didot}, which he fixed at a
little more than the English point, gave the founders a common unit, so that a
line of type cast in one shop would stand true against a line cast in another.

It sounds like a small thing. It was not. A system of sizes that everyone could
rely on meant that a book could be set from the founts of three different
foundries without the compositor having to pack every line with slips of card,
and it meant, too, that the type itself could be made finer, because nothing had
to be thickened to survive a careless fit.

The son, Firmin, took the next step. His roman, cut in the 1780s and refined for
twenty years afterwards, pushed the contrast between thick and thin strokes
further than anyone had dared: the hairlines are almost threads, the serifs flat
and unbracketed, the stress upright, so that a capital \emph{O} looks less like a
written circle than like a drawing made with a ruling pen. Readers who were used
to the warmer letters of Garamond or Caslon found it cold, and some said so in
print. Others found it exactly right for an age that admired order and clear
outlines, and the great editions of Virgil, Horace and Racine that came from the
Didot press between 1798 and 1801 were bought, admired and imitated across
Europe.

\textbf{The bold} of a Didot is a nineteenth-century idea: the founders who
followed Didot and Bodoni thickened the stems and left the hairlines thin, so
that a heading in \textbf{bold} keeps the family's contrast, and a word set in
\textbf{\emph{bold italic}} still reads as part of the same page.
\end{multicols}

\end{document}
